% ------------------------------------------------------------------------------
% Formatvorlage f�r Masterarbeiten an der Ohm-Hochschule N�rnberg
% ------------------------------------------------------------------------------
%   erstellt von Stefan Macke, 24.04.2009
%   http://blog.stefan-macke.de

% Dokumentenkopf ---------------------------------------------------------------
%   Diese Vorlage basiert auf "scrreprt" aus dem koma-script.
% ------------------------------------------------------------------------------
\documentclass[
    11pt, % Schriftgr��e
    DIV10,
    ngerman, % f�r Umlaute, Silbentrennung etc.
    a4paper, % Papierformat
    oneside, % einseitiges Dokument
    titlepage, % es wird eine Titelseite verwendet
    parskip=half, % Abstand zwischen Abs�tzen (halbe Zeile)
    headings=normal, % Gr��e der �berschriften verkleinern
    listof=totocnumbered, % Verzeichnisse im Inhaltsverzeichnis auff�hren
    bibliography=totocnumbered, % Literaturverzeichnis im Inhaltsverzeichnis auff�hren
    index=totoc, % Index im Inhaltsverzeichnis auff�hren
    captions=tableheading, % Beschriftung von Tabellen unterhalb ausgeben
    final % Status des Dokuments (final/draft)
]{scrreprt}

% Meta-Informationen -----------------------------------------------------------
%   Informationen �ber das Dokument, wie z.B. Titel, Autor, Matrikelnr. etc
%   werden in der Datei Meta.tex definiert und k�nnen danach global
%   verwendet werden.
% ------------------------------------------------------------------------------
% Meta-Informationen -----------------------------------------------------------
%   Definition von globalen Parametern, die im gesamten Dokument verwendet
%   werden k�nnen (z.B auf dem Deckblatt etc.).
%
%   ACHTUNG: Wenn die Texte Umlaute oder ein Esszet enthalten, muss der folgende
%            Befehl bereits an dieser Stelle aktiviert werden:
%            \usepackage[latin1]{inputenc}
% ------------------------------------------------------------------------------
\newcommand{\school}{Technische Hochschule N�rnberg Georg Simon Ohm}
\newcommand{\titel}{Assistenzplaner}
\newcommand{\untertitel}{Serverbasierte Software zur Koordinierung der Assistenten f�r behinderte Menschen}
\newcommand{\art}{Masterarbeit}
\newcommand{\studiengang}{Software Engineering und Informationstechnik}
\newcommand{\autor}{Ulrich Belitz}
\newcommand{\fakultaet}{ Elektrotechnik Feinwerktechnik Informationstechnik}
\newcommand{\matrikelnr}{2464776}
\newcommand{\erstgutachter}{Prof. Dr. Thomas Mahr}
\newcommand{\zweitgutachter}{ToDo} %TODO Zweigutachter
\newcommand{\jahr}{2014}
\newcommand{\ort}{Taunusstein}
\newcommand{\logo}{LogoOhmHochschule.jpg} %TODO Logo-Update


% ben�tigte Packages -----------------------------------------------------------
%   LaTeX-Packages, die ben�tigt werden, sind in die Datei Packages.tex
%   "ausgelagert", um diese Vorlage m�glichst �bersichtlich zu halten.
% ------------------------------------------------------------------------------
\input{Packages}

% Erstellung eines Index und Abk�rzungsverzeichnisses aktivieren ---------------
\makeindex
\makenomenclature

% Kopf- und Fu�zeilen, Seitenr�nder etc. ---------------------------------------
\input{Seitenstil}

% eigene Definitionen f�r Silbentrennung
\include{Silbentrennung}

% eigene LaTeX-Befehle
% Eigene Befehle und typographische Auszeichnungen f�r diese

% einfaches Wechseln der Schrift, z.B.: \changefont{cmss}{sbc}{n}
\newcommand{\changefont}[3]{\fontfamily{#1} \fontseries{#2} \fontshape{#3} \selectfont}

% Abk�rzungen mit korrektem Leerraum 
\newcommand{\ua}{\mbox{u.\,a.\ }}
\newcommand{\zB}{\mbox{z.\,B.\ }}
\newcommand{\dahe}{\mbox{d.\,h.\ }}
\newcommand{\Vgl}{Vgl.\ }
\newcommand{\bzw}{bzw.\ }
\newcommand{\evtl}{evtl.\ }

\newcommand{\abbildung}[1]{Abbildung~\ref{fig:#1}}
\newcommand{\gleichung}[1]{Gleichung~\ref{eq:#1}}
\newcommand{\tabelle}[1]{Tabelle~\ref{tab:#1}}
\newcommand{\codelisting}[1]{Listing~\ref{lst:#1}}

\newcommand{\bs}{$\backslash$}

% erzeugt ein Listenelement mit fetter �berschrift 
\newcommand{\itemd}[2]{\item{\textbf{#1}}\\{#2}}

% einige Befehle zum Zitieren --------------------------------------------------
\newcommand{\Zitat}[2][\empty]{\ifthenelse{\equal{#1}{\empty}}{\citep{#2}}{\citep[#1]{#2}}}

% zum Ausgeben von Autoren
\newcommand{\AutorName}[1]{\textsc{#1}}
\newcommand{\Autor}[1]{\AutorName{\citeauthor{#1}}}

% Links mit Zuletzt angerufen am in Fu�note---
%\Link{URL}{Beschreibung}{Datum}
\newcommand{\Link}[3]{\href{#1}{#2}\footnote{\url{#1} - zuletzt abgerufen am #3}}
\newcommand{\LinkWithoutDate}[2]{\href{#1}{#2}\footnote{\url{#1}}}

% neue Abkuerzung - z. B. \NewAbb{UML}{Unified Modling Language}
\newcommand{\NewAbb}[2]{\textit{#2} (\textbf{#1})\nomenclature{#1}{#2}}

% Notizen
\newcommand{\Note}[1]{\par\noindent\colorbox{hellrot}{\parbox{\dimexpr\textwidth-2\fboxsep\relax}{#1}}}

% verschiedene Befehle um W�rter semantisch auszuzeichnen ----------------------
\newcommand{\NeuerBegriff}[1]{(\textbf{#1})}
\newcommand{\Fachbegriff}[1]{\textit{#1}}

\usepackage{tikz}

\newcommand{\foldericon}{%
\begin{tikzpicture}[inner sep=0pt]
 \node (base)[rectangle,minimum width=.8em,minimum height=.8em,draw,rounded
corners=1pt,fill=yellow!50] {};
 \node (paper)[rectangle,minimum width=.7em,minimum height=.7em,draw,rounded
corners=1pt,fill=white] {};
 \node (slant) [rectangle,minimum width=.8em,minimum height=.8em,draw,rounded
corners=1pt,fill=yellow!50,xslant=-.5pt,anchor=south west] at (base.south west)
{};
\end{tikzpicture}
}

\newcommand{\fileicon}{%
\begin{tikzpicture}[inner sep=0pt]
 \node (base)[rectangle,minimum width=.8em,minimum height=.8em,draw,rounded
corners=1pt,fill=white!50] {};
 \node (line2)[rectangle,minimum width=.6em,minimum height=.02em,draw,fill=white!50] {};
 \node (line1)[rectangle,minimum width=.6em,minimum height=.02em,draw,fill=white!50,yshift=.2em] {};
 \node (line3)[rectangle,minimum width=.6em,minimum height=.02em,draw,fill=white!50,yshift=-.2em] {};
%  \node (line1)[rectangle,minimum width=.8em,minimum height=.09em,draw,fill=white!50,yslant=-.5pt,anchor=south west] at (base.south west){};
 %\node (paper)[rectangle,minimum width=.8em,minimum height=.8em,draw,rounded
%corners=1pt,fill=white] {};
 %\node (slant) [rectangle,minimum width=.9em,minimum height=.9em,draw,rounded
%corners=1pt,fill=yellow!50,xslant=-.5pt,anchor=south west] at (base.south west)
%{};
\end{tikzpicture}
}

\newcommand{\file}[1]{\mbox{\fileicon{}\texttt{#1}}}
\newcommand{\folder}[1]{\mbox{\foldericon{}\texttt{#1}}}
\newcommand{\button}[1]{Button \enquote{#1}}

\newcommand{\Eingabe}[1]{\texttt{#1}}
\newcommand{\Code}[1]{\lstinline[language=Java, showstringspaces=false,basicstyle=\ttfamily\color{black}\small,tabsize=2]$#1$}
\newcommand{\CodeX}[1]{\lstinline[language=XML, showstringspaces=false,basicstyle=\ttfamily\color{black}\small,tabsize=2]$#1$}

\newcommand{\Datei}[1]{\texttt{#1}}

\newcommand{\Datentyp}[1]{\textsf{#1}}
\newcommand{\XMLElement}[1]{\textsf{#1}}
\newcommand{\Webservice}[1]{\textsf{#1}}

% Das eigentliche Dokument -----------------------------------------------------
%   Der eigentliche Inhalt des Dokuments beginnt hier. Die einzelnen Seiten
%   und Kapitel werden in eigene Dateien ausgelagert und hier nur inkludiert.
% ------------------------------------------------------------------------------
\begin{document}
\pagenumbering{arabic}
% auch subsubsection nummerieren
\setcounter{secnumdepth}{1}
\setcounter{tocdepth}{1}

% Deckblatt und Abstract ohne Seitenzahl
%\ofoot{}
\thispagestyle{plain}
\begin{titlepage}
\begin{center}

\huge{\textbf{\titel}}\\[1.5ex]
\LARGE{\textbf{\untertitel}}\\[24ex]

\LARGE{\textbf{\art}}\\[1.5ex]
\Large{im Studiengang \studiengang}\\[6ex]

\includegraphics[scale=0.2]{../Images/Thesis/LogoOhmHochschuleMitText.png}\\[6ex]

\normalsize
\begin{tabular}{w{3cm}p{9cm}}\\
vorgelegt von:  & \autor\\[1.2ex]
Fakult�t: & \fakultaet\\[1.2ex]
Matrikelnummer: & \matrikelnr\\[1.2ex]
Erstgutachter:  & \erstgutachter\\[1.2ex]
Zweitgutachter: & \zweitgutachter\\[3ex]
Schlagw�rter: & \keywords
\end{tabular}
\end{center}
\end{titlepage}
\include{Erklaerung} % Selbst�ndigkeitserkl�rung

\addchap{Abstract}
\label{sec:Abstract}
%TODO Abstract schreiben

\addchap{Kurzfassung}
\label{sec:Zusammenfassung}
%TODO Zusammenfassung schreiben

%\ofoot{\pagemark}

% Seitennummerierung -----------------------------------------------------------
%   Vor dem Hauptteil werden die Seiten in gro�en r�mischen Ziffern 
%   nummeriert.
% ------------------------------------------------------------------------------
\tableofcontents % Inhaltsverzeichnis

% Abk�rzungsverzeichnis --------------------------------------------------------

%\listoffigures % Abbildungsverzeichnis



% arabische Seitenzahlen im Hauptteil ------------------------------------------
%\clearpage
%\include{Inhalt/Roadmap}

INHALT
%TODO Inhalt einf�gen


\listoftables % Tabellenverzeichnis
\listoffigures
\renewcommand{\lstlistlistingname}{Verzeichnis der Listings}
\lstlistoflistings % Listings-Verzeichnis

%  Bitte neue Begriffe in alphabetischer Reihenfolge hinzufügen
\nomenclature{EFI}{Fakultät Elektrotechnik Feinwerktechnik}
\nomenclature{HTML}{ToDo}
\nomenclature{HTTP}{HyperText Transport Protocol}
\nomenclature{UML}{Unified Modeling Language}
\nomenclature{URL}{Uniform Resource Locator}
\let\stdchapter\chapter %Nummer vor Abk�rzungsverzeichnis (speichern der chapter Umgebung)
\def\chapter*#1{\stdchapter{#1}} % neue Definition von chapter
\printnomenclature
\label{sec:Glossar}
\let\chapter\stdchapter % zur�ck zur alten Definition

% Literaturverzeichnis ---------------------------------------------------------
%   Das Literaturverzeichnis wird aus der BibTeX-Datenbank "Bibliographie.bib"
%   erstellt.
% ------------------------------------------------------------------------------

\bibliography{../BibTeX/literatur} % Aufruf: bibtex Masterarbeit
\bibliographystyle{natdin} % DIN-Stil des Literaturverzeichnisses

% Anhang -----------------------------------------------------------------------
%   Die Inhalte des Anhangs werden analog zu den Kapiteln inkludiert.
%   Dies geschieht in der Datei "Anhang.tex".
% ------------------------------------------------------------------------------
\begin{appendix}
    \clearpage
    \chapter{Anhang}
    \label{sec:Anhang}
    % Rand der Aufz�hlungen in Tabellen anpassen
    \setdefaultleftmargin{1em}{}{}{}{}{}

    
    
%    \section{Quellcode}

\subsection{OhmSens}
\lstset{language=Java, basicstyle=\ttfamily\color{black}\scriptsize, showstringspaces=false,tabsize=2}

%F�r die Listings wurde ein symbolischer link (Listings/source) in das Projektverzeichnis (Bewegungsanalyse\Sensor\Code\OhmSens) erstellt.

\subsubsection{\Code{package de.ohmhochschule.efi.haenel.ohmsens}}
\subsubsection{\Datei{Main.java}}
\lstinputlisting[label=lst:Main.java,caption=Main.java]{Listings/source/src/de/ohmhochschule/efi/haenel/ohmsens/Main.java}

\subsubsection{\Datei{Prefs.java}}
\lstinputlisting[label=lst:Prefs.java,caption=Prefs.java]{Listings/source/src/de/ohmhochschule/efi/haenel/ohmsens/Prefs.java}

\subsubsection{\Code{package de.ohmhochschule.efi.haenel.ohmsens.service}}

\subsubsection{\Datei{ClassificationDataSet.java}}
\lstinputlisting[label=lst:ClassificationDataSet.java,caption=ClassificationDataSet.java]{Listings/source/src/de/ohmhochschule/efi/haenel/ohmsens/service/ClassificationDataSet.java}

\subsubsection{\Datei{DataSet.java}}
\lstinputlisting[label=lst:DataSet.java,caption=DataSet.java]{Listings/source/src/de/ohmhochschule/efi/haenel/ohmsens/service/DataSet.java}

\subsubsection{\Datei{Distributor.java}}
\lstinputlisting[label=lst:Distributor.java,caption=Distributor.java]{Listings/source/src/de/ohmhochschule/efi/haenel/ohmsens/service/Distributor.java}

\subsubsection{\Datei{FeatureThread.java}}
\lstinputlisting[label=lst:FeatureThread.java,caption=FeatureThread.java]{Listings/source/src/de/ohmhochschule/efi/haenel/ohmsens/service/FeatureThread.java}

\subsubsection{\Datei{FeatureVectorDataSet.java}}
\lstinputlisting[label=lst:FeatureVectorDataSet.java,caption=FeatureVectorDataSet.java]{Listings/source/src/de/ohmhochschule/efi/haenel/ohmsens/service/FeatureVectorDataSet.java}

\subsubsection{\Datei{HandleDataSetThread.java}}
\lstinputlisting[label=lst:HandleDataSetThread.java,caption=HandleDataSetThread.java]{Listings/source/src/de/ohmhochschule/efi/haenel/ohmsens/service/HandleDataSetThread.java}

\subsubsection{\Datei{MarkerDataSet.java}}
\lstinputlisting[label=lst:MarkerDataSet.java,caption=MarkerDataSet.java]{Listings/source/src/de/ohmhochschule/efi/haenel/ohmsens/service/MarkerDataSet.java}

\subsubsection{\Datei{NeighbourClassificatorThread.java}}
\lstinputlisting[label=lst:NeighbourClassificatorThread.java,caption=NeighbourClassificatorThread.java]{Listings/source/src/de/ohmhochschule/efi/haenel/ohmsens/service/NeighbourClassificatorThread.java}

\subsubsection{\Datei{OhmsensRemoteService.java}}
\lstinputlisting[label=lst:OhmsensRemoteService.java,caption=OhmsensRemoteService.java]{Listings/source/src/de/ohmhochschule/efi/haenel/ohmsens/service/OhmsensRemoteService.java}

\subsubsection{\Datei{SensorDataSet.java}}
\lstinputlisting[label=lst:SensorDataSet.java,caption=SensorDataSet.java]{Listings/source/src/de/ohmhochschule/efi/haenel/ohmsens/service/SensorDataSet.java}

\subsubsection{\Datei{WriteToFile.java}}
\lstinputlisting[label=lst:WriteToFile.java,caption=WriteToFile.java]{Listings/source/src/de/ohmhochschule/efi/haenel/ohmsens/service/WriteToFile.java}

\subsubsection{\Datei{WriteToSocket.java}}
\lstinputlisting[label=lst:WriteToSocket.java,caption=WriteToSocket.java]{Listings/source/src/de/ohmhochschule/efi/haenel/ohmsens/service/WriteToSocket.java}

\subsubsection{\Datei{IOhmsensRemoteService.aidl}}
\lstinputlisting[label=lst:IOhmsensRemoteService.aidl,caption=IOhmsensRemoteService.aidl]{Listings/source/src/de/ohmhochschule/efi/haenel/ohmsens/service/IOhmsensRemoteService.aidl}

\subsubsection{\Datei{IServiceCallback.aidl}}
\lstinputlisting[label=lst:IServiceCallback.aidl,caption=IServiceCallback.aidl]{Listings/source/src/de/ohmhochschule/efi/haenel/ohmsens/service/IServiceCallback.aidl}

\subsubsection{Ressourcen}
\lstset{language=XML, basicstyle=\ttfamily\color{black}\scriptsize, showstringspaces=false,tabsize=2}

\subsubsection{Layout / layout}
\subsubsection{\Datei{main.xml}}
\lstinputlisting[label=lst:main.xml,caption=main.xml]{Listings/source/res/layout/main.xml}

\subsubsection{\Datei{settings.xml}}
\lstinputlisting[label=lst:settings.xml,caption=settings.xml]{Listings/source/res/layout/settings.xml}

\subsubsection{Men� / menu}
\subsubsection{\Datei{options\_menu.xml}}
\lstinputlisting[label=lst:optionsmenu.xml,caption=options\_menu.xml]{Listings/source/res/menu/options_menu.xml}

\subsubsection{Werte / values}
\subsubsection{\Datei{arrays.xml}}
\lstinputlisting[label=lst:arrays.xml,caption=arrays.xml]{Listings/source/res/values/arrays.xml}

\subsubsection{\Datei{strings.xml}}
\lstinputlisting[label=lst:strings.xml,caption=strings.xml]{Listings/source/res/values/strings.xml}

\subsubsection{XML / xml}
\subsubsection{\Datei{preference.xml}}
\lstinputlisting[label=lst:preference.xml,caption=preference.xml]{Listings/source/res/xml/preference.xml}

\subsubsection{Manifest}
\subsubsection{\Datei{AndroidManifest.xml}}
\lstinputlisting[label=lst:AndroidManifest.xml,caption=AndroidManifest.xml]{Listings/source/AndroidManifest.xml}

\subsection{OhmSensServer}

%\subsection{MATLAB\textsuperscript{\textregistered} Skripte}
%\subsubsection{\Datei{Klassifikation.m}}
%\lstset{language=XML, basicstyle=\ttfamily\color{black}\scriptsize, showstringspaces=false,tabsize=2}
%\definecolor{dkgreen}{rgb}{0,0.6,0}
%\lstset{language=Matlab,
%   keywords={break,case,catch,continue,else,elseif,end,for,function,
%      global,if,otherwise,persistent,return,switch,try,while},
%   basicstyle=\ttfamily\color{black}\scriptsize,
%   keywordstyle=\color{blue},
%   commentstyle=\color{red},
%   stringstyle=\color{dkgreen},
%   numbers=left,
%   numberstyle=\tiny\color{gray},
%   stepnumber=1,
%   numbersep=10pt,
%   backgroundcolor=\color{white},
%   tabsize=2,
%   showspaces=false,
%   showstringspaces=false}
%\lstinputlisting[label=lst:Klassifikation.m,caption=Klassifikation.m]{Listings/Klassifikation.m}

\end{appendix}

% Index ------------------------------------------------------------------------
%   Zum Erstellen eines Index, die folgende Zeile auskommentieren.
% ------------------------------------------------------------------------------
%\printindex

\end{document}
